\section{Benchmark settings}
\label{sec:settings}

In this section we summarise the settings and inputs used for this benchmark comparison of MC event generators for muon and neutrino DIS at the LHC. 
%
We specify  the kinematics, the fiducial regions, the common physics inputs and the observables.
%
Technical implementation details, including the per-generator cards, 
normalisation conventions, and the validation procedure, are collected in full in the public
repository of this project~\cite{benchmarkrepo}, see also App.~\ref{app:implementation}.

\paragraph{Kinematics and fiducial region.}
%
We consider a lepton (neutrino or muon) with energy $E_\ell$ incident on a tungsten (W) nucleus at rest in two channels:
neutral-current scattering 
\be 
\mu^\pm  + A \to \mu^\pm + X \, ,
\ee 
via $\gamma^*/Z$ exchange, and
charged-current scattering
\be
\nu_\ell~(\bar{\nu}_\ell) + A \to \ell^-~(\ell^+) +  X \, , \quad \ell=e,\mu\, ,
\ee
via $W^+~(W^-)$ exchange.  
%
We consider a range of energies between $E_\ell=400$~GeV and $E_\ell=4$~TeV, with representative results for differential distributions shown for $E_\ell=1$ TeV. 
%
Cross-sections are given per nucleon: each generator is run separately on a proton and on a neutron target, and the two are combined as $\sigma = (74\,\sigma_p + 110\,\sigma_n)/184$, as appropriate for a tungsten nucleus and neglecting nuclear effects (which are studied separately in App.~\ref{app:npdf}).
%
Event generation is performed either in the centre-of-mass frame or in the target rest frame, depending on which mode is more computationally efficient.
%
All observables and selection requirements are evaluated in the target rest frame.

The fiducial region of the benchmark is defined by
\begin{equation}
  Q^2 > 4~{\rm GeV}^2 , \qquad W > 3~{\rm GeV},
  \label{eq:fiducial}
\end{equation}
with $Q^2$ and $W^2$ being the standard DIS invariants, $Q^2$ built from
the beam and scattered-lepton momenta and $W^2 = M_N^2 + 2P\cdot q - Q^2$ also
from the momentum $P$ of the target nucleon at rest, with $M_N$ the nucleon mass. 
%
The generators are run with generation cuts looser than, or at most equal to, those of Eq.~\eqref{eq:fiducial}, and the fiducial region is then imposed on the generated final-state events.
%
The lower limit on $Q^2$ keeps the calculation in the region of validity of perturbative QCD, while the lower bound on the invariant mass of the hadronic final state $W$ avoids the region where  hadronisation models are unreliable.
%
For the specific case of charm-production observables the lower limit on $W$ is raised to 5~GeV.

\paragraph{FASER selections.}
\label{sec:faser-selections}
%
The fiducial region of Eq.~\eqref{eq:fiducial} enables a clean comparison of the predictions of event generators, both inclusive and differential, with the corresponding analytic structure-function calculations.
%
Additional selection requirements need to be imposed to match the acceptance of the FASER experiment.
%
For this, we define three further
selections that emulate FASER analyses, all nested 
inside Eq.~\eqref{eq:fiducial}.
%
Different selections are needed since FASER can perform measurements with different detector subcomponents, each with their own set of requirements on the detected final-state particles.  

\paragraph{\it Spectrometer selection.}  The FASER electronic-detector analyses~\cite{FASER:2023zcr,FASER:2024ref,FASER:2026enu2d} select $\nu_\mu$ and $\bar{\nu}_\mu$ charged-current interactions in a $200$~mm diameter fiducial volume of the tungsten target by requiring a single track passing through all three spectrometer stations with momentum $p_\mu > 100$~GeV, lying within $95$~mm of the detector axis and at an angle $\theta < 25$~mrad to it. 
%
Therefore our truth-level proxy for the FASER spectrometer selection is given by
\begin{equation}
  E_{\mu'} > 100~{\rm GeV} , \qquad \theta_{\mu'} < 25~{\rm mrad} \, ,
  \label{eq:tierS}
\end{equation}
where with $\mu'$ we indicate the final-state muon. 
%
Since $\theta_{\mu'} \simeq \sqrt{Q^2(1-y)}/E_{\mu'}$ in this regime, the additional angular requirement on the outgoing muon removes part of the  high-$Q^2$, high-$y$ corner of the fiducial region.


\paragraph{\it Emulsion selection.}  
%
The FASER$\nu$ emulsion neutrino analyses~\cite{FASER:2024hoe,FASER:2026nu680} are vertex-based.
%
The current FASER$\nu$ event selection therefore requires more than four tracks converging on the vertex, of which more than three have $\tan\theta \le 0.1$, and identifies the CC muon by a track penetrating more than one hundred tungsten plates with $p_\mu > 200$~GeV and $\tan\theta > 0.005$, back-to-back with the remaining vertex tracks in azimuth.
%
Taking into account these considerations, 
our truth-level proxy for the FASER$\nu$ emulsion selection is 
\begin{equation}
  E_{\mu'} > 200~{\rm GeV} , \quad
  \theta_{\mu'} > 5~{\rm mrad} , \quad
  N_{0.5} \ge 5 , \quad N_{0.1} \ge 4 , \quad
  \Delta\phi(\mu', X) > \pi/2 ,
  \label{eq:tierE}
\end{equation}
where $N_{0.5}$ and $N_{0.1}$ count charged particles satisfying $\tan\theta < 0.5$ and $<0.1$ respectively, defined at hadron level under the same $c\tau > 10$~mm stability convention as every other observable in this work,\footnote{In the specific case of charmed mesons, the emulsion detectors can reconstruct $\gamma c\tau \sim 3$~mm tracks, and hence a $D^\pm$ appears as one track and a $D^0$ as none. } and $\Delta\phi(\mu', X)$ is the azimuthal separation between the scattered lepton and the vector sum of all stable charged hadrons ($X=\sum h^\pm$).  
%
In addition, we impose that reconstructed charged particles have energies $E_{\rm ch}\ge 1$~GeV, given that tracks need to pass several tungsten plates in order to be reconstructed.
%
This cut also improves the infrared robustness of the FASER observables, since soft particles, such as the MeV-energy fragments from the breakup of the nuclear remnant in real events, are excluded.

The efficiency of the emulsion selection depends sensitively 
on the charged-particle multiplicity, one of the observables for which, as we will show in Sect.~\ref{sec:hadronic}, event generators differ the most. 
%
The same sensitivity propagates into the
experimental estimate for the incoming neutrino energy $E_\nu$.
%
For instance,  the $2026$ emulsion analysis reconstructs $E_\nu$ with a multivariate estimator whose inputs include the summed
charged-hadron momenta~\cite{FASER:2026nu680}, and hence the modelling of hadronic activity directly affects
the unfolded cross-section estimate. 
%
Another interesting observable in this context is the fraction of energy going into (unobserved) neutral hadrons, which affects the estimate of the hadronic energy resolution.

\paragraph{\it Dimuon selection.} 
%
We also consider an extension of the spectrometer selection where we request a secondary  muon of charge opposite to the primary, the so-called dimuon selection:
\begin{equation}
  \textrm{Spectrometer~selection} \quad \textrm{and} \quad \exists\, \mu^{\pm} :\;
  E_{\mu^{\pm}} > 100~{\rm GeV} , \quad \theta_{\mu^{\pm}} < 25~{\rm mrad} \, ,
  \label{eq:tierD}
\end{equation}
where a secondary $\mu^+$ is associated with a primary $\mu^-$ and vice versa.
%
This selection is not a further acceptance requirement but rather a charm tag, given that an
opposite-sign secondary muon with $E_\mu > 100$~GeV comes overwhelmingly from the
semileptonic decay of a charm hadron.
%
Its rate is set by the charm fraction, a semileptonic branching ratio of $\mathcal{O}(10\%)$, and the acceptance for the decay muon, and it is of order $10^{-3}$ of the inclusive DIS cross-section.
%
Being a charm tag, the dimuon selection is evaluated with $W > 5$~GeV, as all charm-production observables. 

\paragraph{Physics inputs.}

The common physics specifications to which all generators are configured for the benchmark study, unless otherwise specified, are collected in Table~\ref{tab:settings}.
%
Whenever variations with respect to the parameters listed in this table are implemented, these are mentioned specifically.
%
For instance, as mentioned in 
Sect.~\ref{sec:generators}, not all of these options can be implemented in \genie.
%
Concerning proton structure, we adopt the {\tt NNPDF40\_nnlo\_as\_01180} PDF set~\cite{NNPDF:2021njg} via the \lhapdf interface~\cite{Buckley:2014ana} throughout. 
%
As one of the applications to be discussed in Sect.~\ref{sec:faser-app}, predictions for alternative PDF sets are also provided. 
%
Since the target is a tungsten nucleus, we apply isospin symmetry to construct the neutron PDFs from the proton ones. 
%
Nuclear PDF effects are not part of the benchmark settings and are studied separately in App.~\ref{app:npdf}.
%
The analytic structure-function calculations  of \yadism include target-mass corrections~\cite{Schienbein:2007gr}.

%%%%%%%%%%%%%%%%%%%%%%%%%%%%%%%%%%%%%%%%%%
\begin{table}[!tbp]
  \centering
  \small
  \renewcommand{\arraystretch}{1.5}
  \begin{tabularx}{\textwidth}{Xl}
    \toprule
    Setting & Value \\
    \midrule
    Beams                     & Fixed-$E_\ell$ $\mu^-$ (NC) or $\nu_\mu$ (CC) on a tungsten nucleus at rest \\
    Fiducial region           & $Q^2 > 4$~GeV$^2$, $W > 3$~GeV ($W > 5$~GeV for charm production) \\
    Parton distributions      & {\tt NNPDF40\_nnlo\_as\_01180}~\cite{NNPDF:2021njg} via \lhapdf~\cite{Buckley:2014ana} \\
    Renormalisation and factorisation scales                    & $\mu_R = \mu_F = Q$ \\
    Electroweak scheme        & $G_\mu$-matched: $\alpha_{\rm em} = 1/132.119$, $\sin^2\theta_W = 0.22291$ \\
    QED radiation             & switched off \\
    Heavy-quark scheme        & massless charm and bottom (zero-mass, five-flavour scheme) \\
    Target-mass corrections   & included in the structure-function calculations \\
    Stable-particle criterion & $c\tau > 10$~mm \\
    Multiple interactions     & switched off \\
    Final-state interactions     & switched off \\
    \bottomrule
  \end{tabularx}
  \vspace{0.2cm}
  \caption{The common specification to which the event generators are configured for the benchmark study.
    }
  \label{tab:settings}
\end{table}
%%%%%%%%%%%%%%%%%%%%%%%%%%%%%%%%%%%%%%%%%%

The renormalisation and factorisation scales are set dynamically to $\mu_F=\mu_R=Q$, and their variation using the seven-point prescription~\cite{NNPDF:2019ubu} is used to estimate missing higher-order uncertainties (MHOUs), whose size we quote as half the width of the resulting band.
%
We use the $G_\mu$-matched electroweak scheme and switch off QED radiation in the simulations; the impact of the latter is studied separately in Sect.~\ref{sec:hadronic}.
%
Multiple parton interactions play no role, since only one of the two colliding particles carries partons, while final-state interactions with the cold nuclear medium are neglected and their study is beyond the scope of this paper.

The heavy-quark scheme is strictly massless, both for charm and bottom quarks, so that generators can be compared against zero-mass structure functions without a scheme ambiguity. 
%
We estimate the impact of charm mass effects both from the analytical structure-function side, using the FONLL general-mass scheme~\cite{Forte:2010ta}, up to NNLO for neutral-current scattering and up to NLO for charged-current scattering, for which the massive coefficient functions, although known at NNLO~\cite{Berger:2016inr,Gao:2017kkx,Caola:2026kvl}, are not implemented in \yadism, and from the event generator side using \powhegtwo for charm production in neutrino DIS~\cite{Buonocore:2024pdv}, which accounts for charm mass effects in the $s\to c$ and $d \to c$ transitions.
%
Whenever mass effects are included in the simulations, we indicate so explicitly. 
%
The massive CC calculation of \powhegtwo is benchmarked with its \yadism counterpart in the fixed-flavour-number scheme with $n_f=3$ (FFNS $n_f=3$), finding excellent agreement. 

The stability criterion which defines a final-state hadron in the detector is $c\tau>10$~mm.
%
The chosen criterion implies that we assume that charmed hadrons, whose $c\tau$ is between about $45$ and $310$~$\mu$m, never appear as final-state particles and therefore  they must be counted at the point of production.
%
While at FASER$\nu$ this is not necessarily true ($D$-mesons may be detected before decay), this distinction is immaterial from the point of view of the present study.

\paragraph{Observables.}
%
Three main groups of physical observables are considered in the present work:

\begin{enumerate}
  \item Observables constructed from the DIS kinematics, namely from the incoming and outgoing lepton momenta and the total energy  of the hadronic final state.
  %
  These are the momentum transfer, $Q^2$, the momentum fraction, $x_{\rm Bj}$, the inelasticity, $y$, and the invariant mass of the hadronic final state, $W$ (all Lorentz invariants) as well as the scattered-lepton  energy $E_{\ell'}$ and angle $\theta_{\ell'}$  and the energy transfer $\nu$ (or equivalently the total hadronic energy $E_h$) in the target rest frame.
  %
  These observables can be evaluated both with analytical structure functions as well as at the event generator level, and hence are used to validate the implementation of the latter. 

  \item Observables depending on the final-state hadronic activity, such as the charged multiplicity $N_{\rm ch}$, the energy of the leading charged hadron, and $\Delta\phi$, the azimuthal separation between the scattered lepton and the vector sum of the charged hadrons.  
  %
  These observables probe the shower structure and the hadronisation model, and are directly relevant to event selection at the LHC far-forward detectors. 

  \item Observables constructed from a charm tag (see below), such as DIS observables restricted to charm-tagged events and the 
   multiplicities and energy spectra of $D^\pm$ and $D^0/\bar{D}^0$ mesons, as well as the ratio between charged and neutral $D$-mesons.  
      
\end{enumerate}

\paragraph{Charm tagging.}
\label{sec:charmdef}
%
To connect predictions from the analytical structure-function framework for charm production with those provided by event generators we need a functional definition of charm tagging. 
%
From the event generator side, a charm event is tagged by the presence of at least one charmed hadron produced in the event after parton showering and hadronisation.
%
This requirement hence applies to open charm of any species, its excited
states, charmed baryons, and charmonium.
%
Furthermore, this charmed hadron must be prompt. 
%
We note that this is a hadron-level definition, chosen to mimic charm tagging in experimental FASER analyses and is different from a charm tag 
on the hard process.
%
For instance, in NC DIS, in addition to charm-initiated scattering, $\gamma^*/Z + c \to c$ and photon--gluon fusion, $\gamma^* g \to c\bar{c}$, this definition also adds charm produced by $g \to c\bar{c}$ splitting inside the shower and the charm carried by the beam remnant, which are absent from the corresponding structure function analytical calculations. 

On the structure function side, neutral-current charm production is defined as the charm-flavour contribution to $F_2$, $F_L$ and $xF_3$ following~\cite{Forte:2010ta}, namely the contribution from those diagrams where the charm quark couples to the virtual photon.
%
Diagrams which have charm in the final state but where the light quark couples to the virtual photon are included in the light component of the NC structure functions. 
%
Up to NLO this definition coincides with a tag on the charm quark in the hard process, and from NNLO onwards the two differ.
%
The hadron-level tag used for the event generators differs from both definitions since it also counts the charm produced in shower splittings and during fragmentation, as well as charm carried by the beam remnant. 
%
Nevertheless, as we will show in Sect.~\ref{sec:charm}, good agreement is found for inclusive and differential charm-tagged cross-sections between the generator and analytical predictions for the FASER kinematics. 

Concerning charged-current DIS, charm production is defined~\cite{Ball:2011mu} as all contributions proportional to the CKM matrix elements $V_{cd}$ or $V_{cs}$, that is, by the diagrams where the charm quark couples to the $W^\pm$ boson.
%
In this case, the partonic definition differs from a tag on the charm quark in the hard-process final state already at LO. 
%
For instance, the channel $\nu \bar{c} \to \mu^- \bar{s}$ carries $|V_{cs}|^2$ and so belongs to the charm structure function definition, but there the charm quark is in the initial state and is absent from the hard-process final state. 
%
Nevertheless, charm tagging may still arise via the $c$ present in the beam remnant as required by the vanishing charm number of the nucleon. 

We note that \sherpa's parton shower treats the charm quark as massless and produces $g \to c\bar c$ at
a rate an order of magnitude higher than the massive showers of \pythia and
\herwig, unphysically inflating the $D$-meson rates.
%
To correct for this, charm tags in \sherpa are obtained by requiring a charm quark in the event hard record as well as a charm hadron produced in
the event.

Charm-production observables are evaluated in a fiducial region with a tighter requirement on the hadronic invariant mass, $W > 5$~GeV, applied to both currents, to the dimuon selection, and also to the inclusive cross-sections that normalise charm fractions.
%
In NC scattering, the purely massless ZM calculation populates charm-initiated events down to the $W > 3$~GeV boundary, whereas in a showered event the anticharm quark is left in the beam remnant, so that a hadronic final state can only be formed above a threshold of $W \simeq 4.4$~GeV.
%
Below this threshold \pythia and \herwig cannot hadronise such events and discard them, depleting the charm-tagged cross-section and degrading the agreement with the analytical charm structure-function calculation.

\paragraph{Theoretical uncertainties.}
%
As mentioned above, MHOUs associated with the truncation of the QCD expansion are estimated through the standard seven-point variation of $\mu_R$ and $\mu_F$ about the central choice $\mu_R = \mu_F = Q$.
%
This is done both in \yadism~\cite{Candido:2024rkr,Candido:2023utz} and in \powheg, in the former case exploiting that the scale dependence of the coefficient functions is known analytically 
from the scale independence of the structure functions, and in the latter case using its built-in event reweighting feature.
%
We have verified that MHOUs estimated using these two independent approaches are fully consistent. 

We also estimate the theoretical uncertainties associated with the nucleon PDF through dedicated studies where again both \yadism and \powheg calculations are repurposed for different PDF sets as well as each set's own
uncertainty, which requires no additional calculation.
%
Specifically, in addition to NNPDF4.0, we consider predictions from CT18~\cite{Hou:2019efy}, MSHT20~\cite{Bailey:2020ooq}, ABMP16~\cite{Alekhin:2017kpj}, and ATLASpdf21~\cite{ATLAS:2021vod}, in all cases at NNLO accuracy and taking the strong coupling to be $\alpha_s(m_Z)=0.118$.
%
Furthermore, results are also presented with GRV98LO~\cite{Gluck:1998xa}, which is the PDF set underlying \genie's baseline tune. 


